\pdfvariable trailerid {[<C2892026090200000000000000000000><C2892026090200000000000000000000>]}
\documentclass[10pt]{article}
\usepackage[margin=0.78in]{geometry}
\usepackage{amsmath,amssymb,amsthm,booktabs,array,microtype,hyperref}
\hypersetup{colorlinks=true,linkcolor=blue,urlcolor=blue}
\providecommand{\CRevisionRound}{2}
\newtheorem{theorem}{Theorem}
\newtheorem{lemma}{Lemma}
\newtheorem{proposition}{Proposition}
\newcommand{\Hh}{\mathbb H^2_\kappa}
\newcommand{\R}{\mathbb R}
\newcommand{\D}{\mathrm D}
\newcommand{\SO}{\mathrm{SO}^{+}(2,1)}
\title{Constant Magnetic Flow on the Hyperbolic Plane:\\Orbit Trichotomy, Critical Horocycles, and Exact Periods}
\author{Route-A source-local certificate HCS-C289}
\date{2 September 2026}
\begin{document}
\maketitle

\begin{abstract}
On the complete simply connected surface of curvature $-\kappa^2$ we solve
$\D_t\dot\gamma=bJ\dot\gamma$ at fixed speed $v$.
For $v>0$, every trajectory is a circle, horocycle, hypercycle, or geodesic,
according as $|b|$ is greater than, equal to, or less than $\kappa v$ (with
$b=0$ separated).  Only the circle chamber closes, with primitive period
$2\pi/\sqrt{b^2-\kappa^2v^2}$.  At equality the Lorentz-frame generator is
nonzero nilpotent, and the orbit is a nonclosed horocycle.  We prove the
classification globally from a raw $3\times3$ moving-frame equation.
\ifnum\CRevisionRound>0
An independent radius--circumference calculation checks the period and fixes
all orientation and normalization conventions.
\fi
\ifnum\CRevisionRound>1
A strict executable certificate reconstructs the matrix independently on 144
rational parameter cells and five boundary cells; it is regression evidence,
not the proof.  Classical ownership and the Route-A rejection are explicit.
\fi
\end{abstract}

\section{Model, conventions, and result}
Let $\Hh$ be the oriented complete simply connected Riemannian surface with
Gaussian curvature $-\kappa^2$, where $\kappa>0$, and let $J$ denote positive
rotation through $\pi/2$.  We study
\begin{equation}\label{eq:magnetic}
 \D_t\dot\gamma=bJ\dot\gamma,\qquad |\dot\gamma|=v,
\end{equation}
for constant $b\in\R$.  Metric compatibility makes the speed constant; we
nevertheless retain $v$ because replacing $b$ by the geodesic curvature is a
frequent source of a missing factor.  Time in every period statement is the
physical parameter $t$ in~\eqref{eq:magnetic}.

The orbit geometry is classical.  Comtet treats the hyperbolic constant-field
problem and its Landau quantization~\cite{Comtet1987}; Adachi treats magnetic
trajectories on constant-holomorphic-curvature spaces~\cite{Adachi1995}.  No
literature originality is claimed for the classification below.  Our purpose
is a convention-complete derivation joined to independently replayable
evidence.

\begin{theorem}[Complete orbit atlas]\label{thm:atlas}
Assume $\kappa>0$ and $v>0$.  Every solution of~\eqref{eq:magnetic} is complete.
Up to an orientation-preserving isometry of $\Hh$ and a time translation:
\begin{enumerate}
 \item if $|b|>\kappa v$, it is a hyperbolic circle and its primitive period is
 \begin{equation}\label{eq:period}
    T=\frac{2\pi}{\sqrt{b^2-\kappa^2v^2}};
 \end{equation}
 \item if $|b|=\kappa v$, it is a nonclosed horocycle;
 \item if $0<|b|<\kappa v$, it is an unbounded hypercycle;
 \item if $b=0$, it is a geodesic.
\end{enumerate}
The sign of $b$ reverses the orientation but not the unoriented curve or its
period.  At the critical face the ambient generator has $A^3=0$ and $A^2\ne0$.
\end{theorem}

\paragraph{Boundary convention.}
If $v=0$, the curve is stationary and the unit-tangent proof below is not
invoked.  The flat case $\kappa=0$ is a limit, not a point in the theorem:
$b\ne0$ gives a Euclidean circle of period $2\pi/|b|$, whereas $b=0$ gives a
straight line.

\section{The constant Frenet system}
For $v>0$ put $T=\dot\gamma/v$ and $N=JT$.  Equation~\eqref{eq:magnetic}
becomes
\begin{equation}\label{eq:frenet}
  \D_tT=bN,\qquad \D_tN=-bT,
\end{equation}
so the signed geodesic curvature with respect to arclength is $k_g=b/v$.
Use the hyperboloid
\[
 \Hh=\{X\in\R^{2,1}:\langle X,X\rangle_L=-\kappa^{-2},\ X_0>0\},
 \qquad \eta=\operatorname{diag}(-1,1,1).
\]
Set $e_0=\kappa X$.  The ambient derivative of a tangent field $Y$ along
$X$ is
\begin{equation}\label{eq:gauss}
 \frac{dY}{dt}=\D_tY+\kappa^2\langle\dot X,Y\rangle X,
\end{equation}
or, in particular, $\dot T=\kappa v e_0+bN$.  With the signs in
\eqref{eq:frenet}, let $F=(e_0,T,N)$ explicitly denote the matrix whose
columns are those three vectors.  It satisfies the right-action equation
$\dot F=FA$, with
\begin{equation}\label{eq:A}
 A=\begin{pmatrix}
 0&\kappa v&0\\
 \kappa v&0&-b\\
 0&b&0
 \end{pmatrix},\qquad A^{\mathsf T}\eta+\eta A=0.
\end{equation}
Changing from row to column frames transposes the displayed ODE but changes
neither the characteristic polynomial nor any conclusion.

\begin{lemma}[Cubic identity]\label{lem:cubic}
For~\eqref{eq:A},
\begin{equation}\label{eq:cubic}
 \det(\lambda I-A)=\lambda(\lambda^2+b^2-\kappa^2v^2),
 \qquad A^3=(\kappa^2v^2-b^2)A.
\end{equation}
At $|b|=\kappa v$, $A^2\ne0$ and $A^3=0$.
\end{lemma}
\begin{proof}
Direct expansion of the three-by-three determinant proves the first identity;
Cayley--Hamilton proves the second.  Substitution $b=\pm\kappa v$ leaves, for
example, the $(1,1)$ entry of $A^2$ equal to $\kappa^2v^2>0$, so the critical
generator is nonzero of nilpotency index three.
\end{proof}

\section{Proof of the classification}
\begin{proof}[Proof of Theorem~\ref{thm:atlas}]
The frame equation has the global solution $F(t)=F(0)e^{tA}$.  Thus every
initial oriented frame gives a solution for all $t$, and transitivity of
$\SO$ on such frames shows that the analysis exhausts all magnetic initial
data.

If $\delta=b^2-\kappa^2v^2>0$,~\eqref{eq:cubic} gives eigenvalues
$0,\pm i\sqrt\delta$.  The one-parameter Lorentz subgroup is elliptic and
fixes a timelike axis.  Intersecting the hyperboloid with the affine plane at
fixed Lorentz distance from that axis gives a hyperbolic circle.  On the
two-dimensional rotating block, put $e_0=(1,0,0)^{\mathsf T}$ in frame
coordinates.  The cubic identity gives
\[
 e^{tA}e_0=e_0+\frac{\sin\theta}{\sqrt\delta}Ae_0
 +\frac{1-\cos\theta}{\delta}A^2e_0,
 \qquad\theta=t\sqrt\delta.
\]
Here $(Ae_0)_T=\kappa v\ne0$, $(A^2e_0)_T=0$, and
$(A^2e_0)_{e_0}=\kappa^2v^2\ne0$.  A position return first forces
$\sin\theta=0$ from the tangent component and then forces
$\cos\theta=1$ from the $e_0$ component.  Thus the base point returns exactly
for $\theta\in2\pi\mathbb Z$, and~\eqref{eq:period} is primitive.

If $\delta=0$, Lemma~\ref{lem:cubic} gives a nontrivial parabolic subgroup.
Its hyperboloid orbits are horocycles.  More specifically,
$e^{tA}e_0=e_0+tAe_0+t^2A^2e_0/2$, whose tangent component is
$\kappa vt$.  It is nonzero for every $t\ne0$, so the base point itself never
returns; the critical horocycle is not closed.

If $\delta<0$, the nonzero eigenvalues are real and the subgroup is
hyperbolic.  Its nongeodesic hyperboloid orbits are hypercycles.  When $b=0$,
the normal is parallel and the curve is the invariant geodesic itself.  The
sign change $b\mapsto-b$ is conjugated by reversing $N$, so it reverses
orientation and preserves the remaining data.
\end{proof}

\ifnum\CRevisionRound>0
\section{Independent geometric recovery of the period}
The matrix proof already establishes the theorem.  A second derivation checks
that the clock and curvature conventions agree.  A circle of hyperbolic
radius $\rho$ on curvature $-\kappa^2$ has unsigned geodesic curvature and
circumference
\begin{equation}\label{eq:circlegeometry}
 |k_g|=\kappa\coth(\kappa\rho),\qquad
 L(\rho)=\frac{2\pi}{\kappa}\sinh(\kappa\rho).
\end{equation}
Because $|k_g|=|b|/v$, a finite $\rho$ exists exactly when
$|b|>\kappa v$.  Moreover
\[
 \sinh^2(\kappa\rho)
 =\frac{1}{\coth^2(\kappa\rho)-1}
 =\frac{\kappa^2v^2}{b^2-\kappa^2v^2}.
\]
Dividing $L(\rho)$ by $v$ reproduces~\eqref{eq:period}.  Since this is an
embedded circle traversed with nonzero constant speed and fixed orientation,
one circumference is its first position return, independently confirming
primitivity.  As
$|b|\downarrow\kappa v$, the radius and the period diverge and the circles
converge, after recentering, to a horocycle.  If $0<|b|<\kappa v$, a curve at
signed distance $r$ from a geodesic has
$|k_g|=\kappa|\tanh(\kappa r)|$; this gives the hypercycle relation
\begin{equation}\label{eq:hypercycle}
 |\tanh(\kappa r)|=\frac{|b|}{\kappa v}<1.
\end{equation}
Thus the geometric and Lorentz classifications agree at every open chamber
and at their common boundary.

\section{What completeness does and does not assert}
Global existence here follows from the constant frame exponential, rather
than from a compactness shortcut (the hyperbolic plane is noncompact).  The
classification is of unparametrized shapes together with their physical-time
parametrization.  It does not identify different field strengths after a
time change.  Nor does elliptic subgroup type assert dynamical stability of
a perturbed nonconstant-field system; no such perturbation theorem is made.

The critical nilpotent statement is deliberately stronger than observing a
triple zero characteristic polynomial: the latter alone would not exclude
$A=0$ or a lower Jordan index.  The explicit check $A^2\ne0$ fixes the
parabolic orbit and its polynomial frame growth.
\fi

\ifnum\CRevisionRound>1
\section{Finite evidence and hostile separation}
The archived JSON contains the Cartesian product
\[
 \kappa,v\in\{\tfrac12,1,2,3\},\qquad
 b\in\{-6,-3,-2,-1,0,1,2,3,6\},
\]
giving 144 orbit cells, plus five boundary cells for zero speed, the two flat
limits, field reversal, and critical nonclosure.  Rational arithmetic fixes
$b^2-\kappa^2v^2$, $b/v$, the shape ratio, and
$(T/2\pi)^2$ wherever a period exists.

The checker is producer-independent: it does not import the producer and
reconstructs~\eqref{eq:A} from each raw triple.  It verifies
$A^{\mathsf T}\eta+\eta A=0$, the cubic identity, full-grid uniqueness, and
critical $A^2\ne0=A^3$.  In every circle cell it also evaluates
$e^{TA}$ at 70 decimal digits and checks the identity monodromy.  A separate
SymPy program rebuilds the characteristic identity.  Fresh-path replay
requires two generated files to be byte-identical to the archive.  A strict
duplicate-rejecting checker faces 43/43 semantic, structural, exact-type, base-return,
exact-YAML, raw-duplicate, and stale-hash attacks; repaired hashes cannot
bless changed mathematics.

These computations test implementation and sign conventions.  They do not
prove Theorem~\ref{thm:atlas} outside the finite grid; the all-parameter proof
is the analytic argument above.

\section{Route-A audit and scope firewall}
The frozen tuple is
\[
 (\mathrm{A0\_FAIL},\mathrm{A1\_WEAK},\mathrm{A2\_FAIL},
  \mathrm{A3\_FAIL},\mathrm{A4\_FORMAL\_HINT}),
\]
with overall verdict \texttt{ROUTE\_A\_REJECTED}.  A1 is only weak because
individual circle trajectories have intrinsic primitive periods, but their
positive-dimensional clean families supply no enumerable isolated primitive
ledger or repetition census.  A magnetic quantum
Hamiltonian is a formal A4 hint, but this package constructs no self-adjoint
domain, operator, or quantum spectrum.  It identifies no arithmetic local data, Euler factors, root numbers,
automorphy, target divisor/counting law, target functional equation, target
zero match, or target Hilbert--P\'olya operator.  Route B is locked false.
The governing literal is \texttt{NO\_BAD\_EULER\_OR\_ROOT\_NUMBER}.

\section{Limitations and declarations}
The theorem assumes a constant field on the simply connected constant-
curvature plane.  Quotients, variable fields, collisions with boundaries,
nonlinear stability, quantum spectral multiplicities, and semiclassical trace
formulae are outside scope.  The flat result is only the stated limit, and the
stationary $v=0$ face is not forced through a unit-tangent formula.

\paragraph{Data and code availability.}
The canonical evidence, producer-independent checker, symbolic reconstruction,
replay, mutation audit, LaTeX source, and release manifest are contained in
this package.  No external dataset is required.

\paragraph{Ethics, conflicts, and funding.}
No human participants, animals, or personal data were involved.  The author
declares no conflict of interest and no external funding.

\paragraph{CRediT and AI-use statement.}
The package author performed conceptualization, formal analysis, software,
validation, writing, and reproducibility curation.  An AI language model
assisted drafting and code generation; every mathematical and executable
claim is exposed to the recorded independent checks, and responsibility
remains with the author.
\fi

\begin{thebibliography}{9}
\bibitem{Comtet1987}
A.~Comtet, \emph{On the Landau Levels on the Hyperbolic Plane}, Annals of
Physics 173 (1987), 185--209.
\href{https://doi.org/10.1016/0003-4916(87)90098-4}{doi:10.1016/0003-4916(87)90098-4}.

\bibitem{Adachi1995}
T.~Adachi, \emph{Kaehler Magnetic Flows for a Manifold of Constant
Holomorphic Sectional Curvature}, Tokyo Journal of Mathematics 18 (1995),
473--483.
\href{https://doi.org/10.3836/tjm/1270043477}{doi:10.3836/tjm/1270043477}.
\end{thebibliography}
\end{document}
